%% Master CV - Niladribihari Mohanta
%% Comprehensive overview of all competencies, experience, and achievements.
%% Use as a master reference when tailoring role-specific CVs.
%%
%% Compile with: lualatex main_example.tex
%% (pdflatex often fails on modern MiKTeX with fontawesome5 font-expansion errors.)

\documentclass[11pt,a4paper,sans]{moderncv}
\moderncvstyle{banking}
\moderncvcolor{blue}

% Force the name and section headings to render in moderncv blue (color1).
% Default banking leaves them black: moderncvstylebanking.sty's \colorlet
% copies (not aliases) the pre-scheme accent colour, so the name colours are
% frozen before \moderncvcolor runs. Re-let them after. \namefont is the hook
% every name-style macro routes through, so this also works on moderncv 2.3.1
% (Debian/Ubuntu apt), which has no \firstnamestyle/\lastnamestyle at all.
\renewcommand*{\namefont}{\fontsize{34}{36}\bfseries\upshape}
\colorlet{firstnamecolor}{color1}
\colorlet{lastnamecolor}{color1}
\colorlet{namecolor}{color1}
\renewcommand*{\sectionstyle}[1]{{\sectionfont\color{color1}#1}}

\usepackage[utf8]{inputenc}
% pdflatex fallback only (the documented engine is lualatex, which skips this
% branch). Without T1 font encoding pdflatex builds accented letters with
% \accent, and the PDF text layer stores them decomposed - `e` + U+0300 rather
% than U+00E8 - so an ATS keyword match on "Genève" fails while the page looks
% right. moderncv 2.5 loads T1 itself under pdflatex; 2.3.1 (Debian/Ubuntu apt)
% does not. \ifpdftex comes from iftex, which every moderncv version loads.
\ifpdftex\usepackage[T1]{fontenc}\fi
% moderncv loads hyperref itself in an \AtEndPreamble hook, so \hypersetup
% must go in an \AtEndPreamble of our own: on moderncv < 2.4 a top-level
% \usepackage{hyperref} clashes with the class's own
% \RequirePackage[unicode]{hyperref}. From 2.4.0 the class passes its options
% through \PassOptionsToPackage instead, which is what removes that clash.
\AtEndPreamble{\hypersetup{
    colorlinks=true,
    linkcolor=blue,
    filecolor=magenta,
    urlcolor=blue,
    pdftitle={Niladribihari Mohanta - CV},
    % Keep pdfpagemode=UseNone: this block runs after moderncv's own
    % \AtEndPreamble (moderncv.cls sets pdfpagemode there), so a FullScreen
    % value here would win and open every CV in fullscreen presentation mode.
    pdfpagemode=UseNone,
}}
\usepackage[scale=0.80]{geometry}
\usepackage{import}
\usepackage{needspace}
% Do NOT wrap \cventry entries in an outer `itemize`. moderncv sizes each
% entry's tabular to \maincolumnwidth (= \textwidth - \leftskip - \rightskip,
% computed once), but inside an itemize the available width is \linewidth,
% narrower by the list indent. The mismatch produces an Overfull \hbox of
% exactly one em (10.95pt at 11pt) on every entry, title line and body alike.
% Calling \cventry at top level, with an itemize only for its bullet body,
% matches how moderncv's own template uses it and compiles clean.

% personal data
\name{Niladribihari}{Mohanta}
% If you have no address to list, DELETE this whole line. \address{}{}{} fails
% with "There's no line here to end" on every moderncv version.
\address{Visakhapatnam, Andhra Pradesh, India}{}{}
\phone[mobile]{+91 637 124 9887}
\email{niladrigudu@gmail.com}
\extrainfo{\href{https://github.com/niladri-gudu}{GitHub}, \href{https://niladri.dev}{niladri.dev}, \href{https://www.linkedin.com/in/niladribihari-mohanta}{LinkedIn}}

\begin{document}

\makecvtitle

% ============================================================
%     PROFILE STATEMENT
% ============================================================

\vspace{6pt}
\small{Backend engineer specialising in distributed systems, asynchronous event-driven pipelines, and LLM-powered applications. Designed production queue workflows with dead-letter handling, replay and idempotent job processing, and owned backend infrastructure end-to-end from self-hosted services and CI/CD through to frontend migration. Works in TypeScript/Node.js with PostgreSQL, Redis and BullMQ, and builds retrieval-augmented generation systems with LangChain and LangGraph.}

% ============================================================
%     CORE COMPETENCIES
% ============================================================

\section{Core Competencies}
\vspace{1pt}
\begin{itemize}

\item \textbf{Distributed Systems}: asynchronous event-driven architecture, message queues and job processing, dead-letter queues, replay mechanics, exponential backoff and idempotency, real-time delivery over Redis Pub/Sub and Socket.IO.

\item \textbf{Backend Engineering}: TypeScript, Node.js, Fastify, Express, NestJS, RESTful API design, authentication workflows, data serialization pipelines, BullMQ, Bull Board queue observability.

\item \textbf{Data \& Storage}: PostgreSQL, Prisma ORM, MongoDB, Redis, cursor-based pagination, pgvector vector embeddings, Supabase, Firebase.

\item \textbf{AI \& LLM Engineering}: retrieval-augmented generation, LangChain, LangGraph agent orchestration, vector search, LLM API integration, prompt-driven pipelines, agentic development with Claude Code.

\item \textbf{Infrastructure \& Delivery}: Docker, Git, CI/CD with Coolify, self-hosted production on Hetzner, GCP and AWS S3, media pipelines with FFmpeg and adaptive HLS.

\end{itemize}

% ============================================================
%     PROFESSIONAL EXPERIENCE
% ============================================================

\section{Professional Experience}
\vspace{3pt}

% --- Most Recent Role ---
\needspace{5\baselineskip}
\cventry{Jul 2026 - Present}{Packaged App Developer}{Accenture}{Ahmedabad, India}{}{\vspace{1pt}
\begin{itemize}
    \item {Completed intensive full-stack training on the MEAN stack (MongoDB, Express.js, Angular, Node.js), focused on scalable enterprise applications.}
    \item {Develops responsive web applications and RESTful APIs with CRUD operations, authentication workflows, and database integration.}
    \item {Works in an Agile environment across code reviews, debugging, testing, and application deployment.}
\end{itemize}}

\vspace{3pt}

% --- Previous Role ---
\needspace{5\baselineskip}
\cventry{Dec 2024 - Jun 2026}{Backend Engineer}{Baatasari}{Visakhapatnam, India}{}{\vspace{1pt}
\begin{itemize}
    \item {Migrated core backend services from Firebase to self-hosted Supabase, eliminating vendor lock-in and reducing monthly infrastructure costs by \textasciitilde25\%.}
    \item {Architected and deployed production infrastructure on Hetzner via Coolify, cutting deployment time by \textasciitilde40\% through automated CI/CD workflows.}
    \item {Re-engineered the web client architecture from Vite to Next.js, reducing initial page load latency by \textasciitilde30\% via SSR caching.}
    \item {Refactored backend endpoints and data serialization pipelines, improving average API response time by \textasciitilde20\% ahead of the GCP migration.}
    \item {Mentored backend developer interns, reviewing PRs and setting technical direction for a small team.}
\end{itemize}}

% ============================================================
%     PROJECTS
% ============================================================

% Keep the heading with the first entry. A section-level \needspace is normally
% avoided because it can strand whitespace and add a page, but here the page
% count stays at 2 and the alternative is an orphaned heading alone at the
% bottom of page 1, which is strictly worse.
\needspace{10\baselineskip}
\section{Projects}
\vspace{3pt}

\needspace{5\baselineskip}
\cventry{Aug 2026 - Sep 2026}{RAG system over live code and memory}{Project Intelligence Agent}{}{}{\vspace{1pt}
\href{https://github.com/niladri-gudu/project-intelligence-agent}{github.com/niladri-gudu/project-intelligence-agent}
\begin{itemize}
    \item {Built an automated repository-intelligence pipeline that clones repos, parses abstract syntax trees, and indexes codebases for contextual retrieval.}
    \item {Designed a RAG architecture with LangChain for chunking, embedding and retrieval, and LangGraph orchestrating stateful multi-step agentic reasoning over codebase context.}
    \item {Implemented asynchronous worker queues for repository diffs, dependency graphing and LLM rate-limit throttling, cutting average query latency by \textasciitilde30\% and improving retrieval relevance by \textasciitilde35\%.}
\end{itemize}}

\vspace{3pt}

\needspace{5\baselineskip}
\cventry{Apr 2026 - May 2026}{Distributed video processing and streaming platform}{TranscodeX}{}{}{\vspace{1pt}
\href{https://github.com/niladri-gudu/transcodex}{github.com/niladri-gudu/transcodex}
\begin{itemize}
    \item {Architected an asynchronous event-driven pipeline using BullMQ and Redis to orchestrate FFmpeg transcoding into 480p/720p/1080p renditions, under default HLS packaging and thumbnail generation.}
    \item {Implemented multipart video chunk uploads and real-time processing telemetry using Redis Pub/Sub and Socket.IO, and integrated automated transcription with contextual search over processed media.}
\end{itemize}}

\vspace{3pt}

\needspace{5\baselineskip}
\cventry{Mar 2026 - Apr 2026}{Distributed social networking backend}{Strata}{}{}{\vspace{1pt}
\href{https://github.com/niladri-gudu/strata}{github.com/niladri-gudu/strata}
\begin{itemize}
    \item {Engineered an event-driven feed and notification engine using BullMQ workers and Redis Pub/Sub, separating ingestion from delivery for consistent performance under load.}
    \item {Designed resilient job-handling workflows with Dead Letter Queues, replay mechanics and exponential backoff retry policies, minimising lost events during failures.}
    \item {Built cursor-based pagination APIs for activity feeds and exposed runtime visibility through Bull Board dashboards, health endpoints and queue metrics.}
\end{itemize}}

% ============================================================
%     EDUCATION
% ============================================================

\section{Education}
\vspace{1pt}

\cventry{2022-2026}{B.Tech in Computer Science and Engineering}{GITAM University}{Visakhapatnam, India}{}{\vspace{1pt}
CGPA: 7.76/10.
}

% ============================================================
%     LANGUAGES
% ============================================================

\section{Languages}
\vspace{1pt}
\begin{itemize}
\item {English (professional, spoken and written), Hindi (conversational, spoken), Odia (conversational, spoken).}
\end{itemize}

% ============================================================
%     REFERENCES
% ============================================================

\section{References}
\vspace{1pt}
\begin{itemize}
\item{Available upon request.}
\end{itemize}
\end{document}
